% ============================================================
% SECCION 4 - MODELOS UML (Danela Arteaga - PE5)
% Fragmento para \input en el informe final
% ============================================================
\section{Modelos UML}
\label{sec:modelos-uml}

El sistema se modeló con el conjunto de diagramas UML disponible en
\texttt{03\_Modelado/Diagramas\_UML/}: un diagrama de casos de uso general (UC-01) con
dieciséis casos de uso especificados (UC-01 a UC-16), doce diagramas de secuencia
(SEQ-01 a SEQ-12) que cubren los flujos de los requisitos funcionales de prioridad
\emph{Debe tener} migrados de la Entrega 2 (1B), un diagrama de clases refinado (CD-01) con
las once clases del dominio (\texttt{User}, \texttt{Administrator}, \texttt{Farmworker},
\texttt{Plot}, \texttt{Crop}, \texttt{Harvest}, \texttt{Disease}, \texttt{Labor},
\texttt{InventoryItem}, \texttt{Report}, \texttt{AIRecommendation}, \texttt{Alert}), un
diagrama de estados (SD-01) para el ciclo de vida de las tareas de campo, tres diagramas de
actividad (AD-01 a AD-03), un diagrama de despliegue (DEP-01) y un diagrama de componentes
(COMP-01). El equipo complementó el análisis con el marco i* (\texttt{istar\_sr},
\texttt{istar\_sd}) para representar la racionalidad estratégica de los actores, en lugar de un
Diagrama de Flujo de Datos formal.

\subsection{Diagrama de clases y su relación con el módulo de IA}
\label{sec:cd01-interpretacion}

Dos relaciones del diagrama de clases (CD-01) son las que sostienen el componente de
inteligencia artificial del sistema. Primero, \texttt{AIRecommendation} y \texttt{Alert} se
asocian con \texttt{Disease} e \texttt{InventoryItem} respectivamente: una recomendación de
manejo (RF-15) y una alerta temprana de enfermedad (RF-26) se generan a partir de un registro de
\texttt{Disease}, mientras que una alerta por disponibilidad crítica (RF-25) se genera a partir de
un \texttt{InventoryItem} por debajo de su umbral mínimo. Esta doble fuente es la razón por la que
el criterio de verificación de RF-25 se corrigió durante la re-inspección de esta práctica
(hallazgo DR-02, \texttt{09\_Cierre\_PE5/registro\_reinspeccion\_PE5.csv}): al provenir de dos
clases distintas, la verificación no puede tratarse como un único conteo, sino como dos
verificaciones separadas por tipo de origen.

Segundo, el diagrama de clases no contiene ninguna clase para representar ingresos económicos
(\texttt{Income} o equivalente). El requisito RF-23 (registro de ingresos), incorporado en la
Entrega 3 (2A) junto con los otros diez requisitos funcionales nuevos, queda así sin una clase de
dominio a la que asociarse. Se reporta como un hallazgo abierto de esta unidad y no como un error
de trazabilidad: la matriz (Sección~\ref{sec:gestion-trazabilidad}) documenta explícitamente esta
ausencia para que se resuelva en el modelado antes del cierre del proyecto.

\subsection{Ciclo de vida de las tareas de campo}
\label{sec:sd01-interpretacion}

El único diagrama de estados del proyecto (SD-01) modela el ciclo \emph{Pending
$\rightarrow$ InProgress $\rightarrow$ Completed} de una tarea de campo (\texttt{Labor}), con
\emph{Blocked} como sub-estado operativo de \emph{InProgress} —no un cuarto valor del enumerado
\texttt{TaskStatus} del diagrama de clases, que conserva solo tres valores—, introducido para
trazabilidad con el mockup de tablero kanban de la Entrega 2. Este diagrama aplica únicamente a
los requisitos que gobiernan el ciclo de una tarea asignada (RF-04, RF-08, RF-09, RF-10); el resto
de los requisitos del sistema no tiene un ciclo de vida propio y así se declara en la matriz de
trazabilidad.
